\section{Removing a fixed size condition without adding variables}
\label{sec:finite-anchor}

The following elementary observation is useful when the number of variables
has priority over both the degree and the size of the coefficients.  It
replaces a condition on the index of a point in a cyclic group by a finite
product of polynomials.  The length of that product can be extremely large.

\begin{proposition}[Excluding finitely many indices]\label{prop:finite-anchor}
Let $K=\mathbb Q(\sqrt2)$, let $E/K$ be an elliptic curve in a Weierstrass
equation, and let $P\in E(K)$ have infinite order.  Fix an integer $M\geq1$.
There is an explicitly computable polynomial
$A_M\in\mathbb Z[X_0,X_1]$ such that, for every nonzero integer $s$,
\[
 A_M(x_0([s]P),x_1([s]P))\ne0
 \quad\Longleftrightarrow\quad |s|>M.
\]
Here $x(Q)=x_0(Q)+\sqrt2\,x_1(Q)$.  The polynomial has degree $2M$ and
requires no additional existential variables.
\end{proposition}
\begin{proof}
For $1\leq n\leq M$, compute integers $a_n,b_n,D_n$ with $D_n>0$ and
\[
 x([n]P)=(a_n+b_n\sqrt2)/D_n.
\]
Such a representation is obtained by the usual exact group law and clearing
the two rational denominators.  Define
\[
 A_M(X_0,X_1)=\prod_{n=1}^M
 \bigl((D_nX_0-a_n)^2-2(D_nX_1-b_n)^2\bigr).
\]
Each factor is $D_n^2$ times the field norm of
$X_0+\sqrt2 X_1-x([n]P)$.  Since the norm of an element of a field vanishes
only at zero, a factor vanishes precisely when the two abscissas agree.
On a Weierstrass elliptic curve, two points with the same abscissa are
negatives of each other.  Thus $x([s]P)=x([n]P)$ holds precisely when
$[s]P=[n]P$ or $[s]P=[-n]P$.  The infinite order of $P$ gives $s=n$ or
$s=-n$.  The leading homogeneous part of the displayed product is nonzero,
so its degree is $2M$.
\end{proof}

An inequality $A_M\ne0$ is a guard, rather than an extra equation with its
own inverse variable.  In a construction which already has a common
denominator guard, one inserts $A_M$ into that product.  In the finite
algebra construction below, one inserts it into the element whose nonzero
components are selected.  In either case it costs no new coordinate.

\begin{corollary}[Conditional replacement in the repository construction]
\label{cor:556-anchor}
Assume the arithmetic soundness argument of the repository, including its
height estimate with a constant $H<2^{4783227}$, and its completeness with
arbitrarily large admissible indices $s$.  Replace the packed semilocal
anchor by the guard $A_M(x(S))\ne0$, where
\[
                    M=2^{9566463}.
\]
The resulting construction still defines $\mathbb Z$.  Applying the
finite algebra construction to it uses $556$ total variables, including the
input.
\end{corollary}
\begin{proof}
The numerical choice satisfies
\[
 (20H)^2<400\,2^{9566454}<2^{9566463}=M.
\]
Proposition~\ref{prop:finite-anchor} therefore supplies the strict size
condition $|s|>(20H)^2$ used in the soundness argument.  Completeness is
preserved by choosing the admissible index larger than this fixed $M$.
The former packed anchor used three coordinates: two retained coordinates
and one coordinate finite over them.  Removing them lowers the number of
retained coordinates from $557$ to $555$.  The finite algebra construction
adds one coordinate to the retained base, giving $556$ total coordinates.
\end{proof}

This corollary is an arithmetic modification of the defining system.
Unlike elimination of finite fibres, it is not asserted to preserve the
projection of an arbitrary system independently of the source arithmetic
theorem.  It also does not make the enormous finite product practical to
expand.  Its integer coefficients are specified exactly by a terminating
indexed construction.  If the contacts or height theorem are changed, the
constant $M$ must be chosen again from the new height estimate; the
numerical value above must not be reused without that check.

